\section{Введение}

\tab Верификация является проверкой соответствия имплементации схемы её спецификации. Имплементация как правило представляет RTL-код (Register Transfer Level), написанный на HDL (hardware description language). Спецификация написана на человеческом языке, поэтому для верификации устройства ее необходимо интерпретировать, то есть переписать в компьютерный код так, чтобы тем или иным способом выразить архитектурное намерение.

В верификации можно выделить три основных вызова. Во-первых, это количество состояний системы. Согласно закону Мура, сложность интегральных схем растет в два раза каждые два года. Вместе с этим растет и сложность их верификации. Более того, предполагается, что сложность верификации в теории может расти экспоненциально по отношению к сложности устройства, что представляет особую сложность при разработке новых микросхем. Эта проблема становится еще более явной, если учесть, что верификация может занимать до 70\% от всего времени разработки. 

Во-вторых, с помощью функциональной верификации мы можем установить только существование ошибки, но не ее отсутствие, то есть отсутствует способ доказать полную корректность устройства. Частично это следствие из предыдущей проблемы: с ростом числа логических элементов становится невозможным проверить все возможные состояния системы за ограниченное время. Также сама спецификация и/или ее интерпретация может иметь множество недостатков: она может быть неоднозначной, неполной или вовсе неправильной. 

Проблему установления корректности системы в какой-то мере решает сбор покрытия. Верификационное покрытие --- это особая метрика, необходимая для оценки проверенного функционала относительно общего требуемого или для оценки отработанной в тестах логики по отношению ко всей существующей. Иначе говоря, покрытие оценивает прогресс и качество верификации. Следует отметить, что стопроцентное покрытие не гарантирует нам полную корректность DUT. Например, в случае ошибки в спецификации или неполного верификационного плана все тесты могут проходить со стопроцентным покрытием, но устройство может содержать ошибки. Таким образом, информативным для нас является только неполное покрытие, так как оно указывает на части модуля, которые точно не проверялись.

В-третьих, верификация должна выявлять ошибки как можно раньше. Установлено, что с каждым последующим этапом производства СнК стоимость исправления ошибки увеличивается. Основные финансовые потери возникают из-за задержки времени выпуска продукта относительно наиболее благоприятного окна возможностей на рынке (market window). Найти причину и устранить ошибку на ранних этапах разработки занимает значительно меньше времени, чем на этапах, более близких к tapeout --- отправке схемы СнК на фабрику для печати. При этом после tapeout исправление ошибки в интегральной схеме требует перезапуска всего цикла изготовления. 

Сейчас существует множество видов верификации: ПЛИС прототипирование, статическая (формальная) верификация, динамическая верификация (симуляция, в т.ч. UVM), трассы(?). При этом из этих методов нельзя выделить наилучший без знания конкретных характеристик проверяемой модели. 
        
Возникает потребность в изучении их применимости, качества и скорости по отношению к сложности и характерным особенностям проверяемой системы. 

В связи с ростом популярности машинного обучения и активным развитием этой сферы, сейчас с каждым годом растет потребность в графических картах общего назначения (GPGPU). Одним из ключевых компонентов графических карт является подсистема памяти, в частности, кэш второго уровня. 
Основной сложностью верификации кэша является крайне большое количество состояний, которое может принять система. В связи с этим встает вопрос не только о качестве верификации, но также и о затрачиваемом времени по отношению к проверяемым состояниям. 
Целью данной работы является изучение методов формальной верификации, симуляции и трасс, исследование их применимости и эффективности в контексте кэша второго уровня GPGPU, а также сравнение этих методов.

